Esta línea de trabajo es el núcleo del proyecto: un clasificador de lesiones cutáneas construido sobre un codificador preentrenado de forma autosupervisada y adaptado por dominio hacia el contexto clínico de interés, sin depender de grandes volúmenes de datos etiquetados.

\subsubsection{Actividades y fases}

El desarrollo de esta línea siguió cuatro actividades encadenadas, correspondientes a las fases 1--3 descritas en la sección~\ref{sec:fases-investigacion} y trazadas en la Tabla~\ref{tab:trazabilidad-fases}: (i) selección y caracterización de los conjuntos HAM10000 y PAD-UFES-20; (ii) preentrenamiento autosupervisado del codificador sobre imágenes sin etiqueta, con adaptación de dominio hacia PAD-UFES-20; (iii) ajuste de un clasificador lineal (regresión logística) sobre representaciones congeladas, usando etiquetas de HAM10000; y (iv) evaluación comparativa frente al codificador público, con selección de variantes restringida al conjunto interno y evaluación externa reservada.

\subsubsection{Herramientas y arquitectura}

El codificador base es DINOv2 ViT-B/14, un transformador visual preentrenado públicamente que se adapta mediante SimSiam, el único algoritmo autosupervisado propio de esta línea. ViT-B/14 aporta la arquitectura; SimSiam aporta el procedimiento de adaptación. Se estudiaron dos variantes de adaptación: A* (sin término de consistencia de tono) y B* (con dicho término, descrito en la sección~\ref{sec:modelamiento-clasificador}). Sobre el codificador adaptado, un clasificador lineal (regresión logística) se ajusta con las representaciones y las etiquetas de HAM10000.

Entorno de cómputo del preentrenamiento: cluster UNABIA vía SLURM, 1 nodo con 4 GPU, 32 CPU y 128\,GB de RAM, presupuesto de 40--80 GPU-hora según la variante del \textit{backbone}; Python 3.11, PyTorch 2.4.1, CUDA 12.1, precisión \texttt{bf16}, tamaño de lote efectivo 128, semilla 42.

\subsubsection{Datos y particiones}

HAM10000 es la fuente de entrenamiento y evaluación del clasificador, particionada por \texttt{lesion\_id} para que ninguna lesión aparezca en más de un subconjunto a la vez. PAD-UFES-20 es el dominio objetivo de adaptación: la partición \texttt{target\_adapt} (1.141 imágenes, sin etiquetas) se incorpora al preentrenamiento autosupervisado para continuar adaptando el codificador al dominio clínico; la partición \texttt{target\_test} se reserva íntegramente para evaluación externa y queda excluida de todo entrenamiento, selección de modelos o hiperparámetros, calibración y ajuste de umbrales.

La Figura~\ref{fig:flujo-adaptacion-dominio} presenta el contraste entre D2 y D3. Ambos parten de la misma inicialización y continúan el aprendizaje autosupervisado durante 2.000 pasos; D2 usa el pool fuente y D3 incorpora además 1.141 imágenes clínicas sin etiquetas de \texttt{target\_adapt}, aproximadamente el 2,7\,\% de su pool. Después se congela cada encoder y se ajustan, por condición, el escalador y la regresión logística con las mismas particiones etiquetadas de HAM10000. La comparación final utiliza únicamente \texttt{target\_test}, separado por paciente, sin realimentar el entrenamiento ni la selección. Los resultados de este piloto se presentan en el Capítulo~\ref{chap:resultados}.

\begin{figure}[!htbp]
\caption{Adaptación autosupervisada de dominio: control fuente D2 y condición D3.}
\label{fig:flujo-adaptacion-dominio}
\centering
\includesvg[width=\textwidth]{images/organizadas/20_adaptacion_dominio}
\par\smallskip
\raggedright\fontsize{11}{13}\selectfont\singlespacing
\par\smallskip{\fontsize{11}{13}\selectfont\singlespacing\raggedright
\noindent\textit{Nota.} Elaboración propia a partir del protocolo del proyecto. DINOv2 y SimSiam se apoyan en \textcite{oquab2024dinov2,chen2021simsiam}; PAD-UFES-20 se describe en \textcite{pacheco2020padufes}. La pérdida dibujada resume el objetivo SimSiam; los términos activos adicionales se especifican en el modelamiento matemático. D3 conserva datos fuente además de los clínicos. La evaluación externa comprende 567 imágenes de 363 pacientes y un mapeo parcial BCC/MEL/NV; no constituye validación clínica en Santander.\par}
\par
\end{figure}

\subsubsection{Modelamiento matemático}
\label{sec:modelamiento-clasificador}

Para una imagen $x_i$ y dos vistas aumentadas $v_{i1}$ y $v_{i2}$, el codificador $f_\theta$, el proyector $g_\phi$ y el predictor $q_\psi$ producen
\begin{equation}
 h_{ij}=f_\theta(v_{ij}),\qquad j\in\{1,2\}.
\end{equation}
El proyector transforma la representación del codificador:
\begin{equation}
 z_{ij}=g_\phi(h_{ij}).
\end{equation}
El predictor procesa la salida del proyector:
\begin{equation}
 p_{ij}=q_\psi(z_{ij}).
\end{equation}
La pérdida SimSiam detiene el gradiente en los objetivos $z$ y usa similitud coseno negativa simétrica:
\begin{equation}
 \mathcal L_{\mathrm{SSL}}=-\frac{1}{2N}\sum_{i=1}^{N}\left[
 \frac{p_{i1}^{\mathsf T}\operatorname{sg}(z_{i2})}{\|p_{i1}\|_2\|\operatorname{sg}(z_{i2})\|_2}
 +\frac{p_{i2}^{\mathsf T}\operatorname{sg}(z_{i1})}{\|p_{i2}\|_2\|\operatorname{sg}(z_{i1})\|_2}\right],
\end{equation}
donde $N$ es el tamaño del lote y $\operatorname{sg}$ detiene el gradiente de esa rama.

La variante B* agrega un término de consistencia de tono: una tercera vista $v_{it}$ reutiliza el recorte de la primera vista sobre una copia con tono modificado, y con $h_{it}=f_\theta(v_{it})$,
\begin{equation}
 \mathcal L_{\mathrm{tono}}=\frac{1}{N}\sum_{i=1}^{N}\left[
 1-\frac{h_{it}^{\mathsf T}\operatorname{sg}(h_{i1})}{\|h_{it}\|_2\|\operatorname{sg}(h_{i1})\|_2}\right].
\end{equation}
El entrenamiento ancla los pesos entrenables $\Theta_{\mathrm{ent}}$ hacia su estado inicial $\theta^0$ (L2SP) y las características hacia el codificador público congelado $f_{\theta^0}$:

\noindent Anclaje de parámetros (L2SP): penaliza la distancia cuadrática entre los pesos entrenables y sus valores iniciales.
\begin{equation}
 \mathcal L_{\mathrm{L2SP}}=\sum_{a\in\Theta_{\mathrm{ent}}}(\theta_a-\theta_a^0)^2.
\end{equation}
\noindent Anclaje de representaciones: compara las características aprendidas con las del codificador público congelado sobre la misma vista.
\begin{equation}
 \mathcal L_{\mathrm{dist}}=\frac{1}{N}\sum_{i=1}^{N}\left[
 1-\frac{h_{i1}^{\mathsf T}f_{\theta^0}(v_{i1})}{\|h_{i1}\|_2\|f_{\theta^0}(v_{i1})\|_2}\right].
\end{equation}
\noindent Pérdida total: combina la pérdida base y los términos complementarios con los pesos de la configuración.
\begin{equation}
 \mathcal L_{\mathrm{total}}=\mathcal L_{\mathrm{base}}+\lambda_t\mathcal L_{\mathrm{tono}}+
 \lambda_a\mathcal L_{\mathrm{L2SP}}+\lambda_d\mathcal L_{\mathrm{dist}}.
\end{equation}

Solo se calculan los términos con peso activo en la configuración de la corrida; la variante A* fija $\lambda_t=0$. Según la evaluación consolidada del proyecto, el término de tono (B* frente a A*) no mostró una diferencia demostrable en ningún conjunto evaluado: el aporte medible proviene de la adaptación SSL anclada en sí misma, no del término de tono. El resultado completo se presenta en el Capítulo~\ref{chap:resultados}.

Sobre las representaciones estandarizadas ($u_{ij}=(h_{ij}-\mu_j)/\sigma_j$, con escala unitaria para dimensiones de varianza cero), el clasificador lineal produce
\begin{equation}
 P(y_i=k\mid x_i)=\frac{\exp(w_k^{\mathsf T}u_i+b_k)}{\sum_{c=1}^{7}\exp(w_c^{\mathsf T}u_i+b_c)},
\end{equation}
La predicción corresponde a la clase con mayor probabilidad:
\begin{equation}
 \widehat y_i=\operatorname*{arg\,max}_{k}P(y_i=k\mid x_i).
\end{equation}

\subsubsection{Algoritmo de preentrenamiento}

El Algoritmo~\ref{alg:simsiam-anclado} resume el flujo del clasificador: adaptación autosupervisada, extracción congelada y ajuste de la sonda supervisada.

\begin{algorithm}[htbp]
\fontsize{11}{13}\selectfont\singlespacing
\caption{Flujo del clasificador SSL: adaptación y clasificación}
\label{alg:simsiam-anclado}
\begin{algorithmic}[1]
\State \textbf{Input:} Conjunto de imágenes sin etiqueta $\{x_i\}_{i=1}^N$ (HAM10000 + PAD-UFES-20 \texttt{target\_adapt}); codificador público preentrenado $f_{\theta^0}$; pesos $\lambda_t,\lambda_a,\lambda_d$; número de pasos $T$
\State \textbf{Initialize:} $\theta\leftarrow\theta^0$ \Comment{copia entrenable del codificador público}
\For{$t=1$ \textbf{to} $T$}
  \State Muestrear lote de imágenes; generar vistas aumentadas $v_{i1},v_{i2}$ (y $v_{it}$ si la variante usa término de tono) \Comment{Aumentación}
  \State Calcular $h_{ij}=f_\theta(v_{ij})$, $z_{ij}=g_\phi(h_{ij})$, $p_{ij}=q_\psi(z_{ij})$ \Comment{Paso hacia adelante}
  \State Calcular $\mathcal L_{\mathrm{SSL}}$ con parada de gradiente en $z$ \Comment{Pérdida SimSiam}
  \If{variante B*}
    \State Calcular $\mathcal L_{\mathrm{tono}}$ con $h_{it}$ \Comment{Consistencia de tono}
  \EndIf
  \State Calcular $\mathcal L_{\mathrm{L2SP}}$ y $\mathcal L_{\mathrm{dist}}$ contra $\theta^0$ y $f_{\theta^0}$ \Comment{Anclaje}
  \State $\mathcal L_{\mathrm{total}}\leftarrow \mathcal L_{\mathrm{SSL}}+\lambda_t\mathcal L_{\mathrm{tono}}+\lambda_a\mathcal L_{\mathrm{L2SP}}+\lambda_d\mathcal L_{\mathrm{dist}}$
  \State Actualizar $\theta,\phi,\psi$ por descenso de gradiente sobre $\mathcal L_{\mathrm{total}}$
\EndFor
\State Congelar $f_\theta$ y extraer representaciones de HAM10000
\State Ajustar el escalador y la regresión logística con la partición de entrenamiento etiquetada
\State Aplicar el preprocesamiento fijado y obtener probabilidades de las siete clases
\State Evaluar en las particiones reservadas; no ajustar con \texttt{target\_test}
\State \Return codificador adaptado, escalador y clasificador
\end{algorithmic}
\end{algorithm}

Sobre $f_\theta$ ya adaptado, un clasificador lineal se ajusta por separado con etiquetas de HAM10000 (regresión logística, ecuación anterior); este paso no forma parte del algoritmo autosupervisado.

\subsubsection{Métricas de evaluación}\label{sec:metricas-clasificacion}

Para $K=7$ clases, la matriz de confusión $C\in\mathbb{N}^{K\times K}$ define los conteos de cada clase $k$: $\mathrm{TP}_k=C_{kk}$, $\mathrm{FP}_k=\sum_{l\neq k}C_{lk}$ y $\mathrm{FN}_k=\sum_{l\neq k}C_{kl}$. Las definiciones generales de exactitud, precisión, exhaustividad y F1 se presentan en la Sección~\ref{sec:metricas-evaluacion}; aquí se especifica su aplicación multiclase y las métricas adicionales del proyecto.

\begin{itemize}[label={\raisebox{0.15ex}{\rule{0.45ex}{0.45ex}}},leftmargin=1.5em,labelsep=0.6em,itemsep=0.4\baselineskip,topsep=0.3\baselineskip]
\item Precisión por clase: proporción de predicciones de la clase $k$ que corresponden a esa clase.
\begin{equation}
 \mathrm{Prec}_k=\frac{\mathrm{TP}_k}{\mathrm{TP}_k+\mathrm{FP}_k}.
\end{equation}

\item Exhaustividad por clase: proporción de casos de la clase $k$ que fueron identificados correctamente.
\begin{equation}
 \mathrm{Rec}_k=\frac{\mathrm{TP}_k}{\mathrm{TP}_k+\mathrm{FN}_k}.
\end{equation}

\item F1 por clase: media armónica de la precisión y la exhaustividad de la clase $k$.
\begin{equation}
 F1_k=\frac{2\,\mathrm{Prec}_k\,\mathrm{Rec}_k}{\mathrm{Prec}_k+\mathrm{Rec}_k}.
\end{equation}

\item F1 macro: promedio no ponderado de F1 entre las clases. Es la métrica primaria del proyecto, pues asigna el mismo peso a cada clase.
\begin{equation}
 F1_{\mathrm{macro}}=\frac{1}{K}\sum_{k=1}^{K}F1_k.
\end{equation}

\item Exactitud balanceada: promedio de la sensibilidad de las clases.
\begin{equation}
 \text{Ex. balanceada}=\frac{1}{K}\sum_{k=1}^{K}\mathrm{Rec}_k.
\end{equation}

\item AUC-ROC por clase: en el enfoque uno-contra-el-resto (OvR), compara las puntuaciones $s_i$ de los $n_k^+$ casos positivos y los $n_k^-$ negativos de la clase $k$, contando los empates con peso de un medio.
\begin{equation}
 \mathrm{AUC}_k=\frac{1}{n_k^{+}n_k^{-}}\sum_{i\in k^{+}}\sum_{j\in k^{-}}\left[\mathbb{1}(s_i>s_j)+\tfrac{1}{2}\mathbb{1}(s_i=s_j)\right].
\end{equation}

\item AUC-ROC macro OvR: promedio no ponderado del AUC de las clases.
\begin{equation}
 \mathrm{AUC}_{\mathrm{OvR}}=\frac{1}{K}\sum_{k=1}^{K}\mathrm{AUC}_k.
\end{equation}

\item Error de calibración esperado (ECE): diferencia entre exactitud y confianza en $M$ contenedores $B_m$, ponderada por su número de observaciones.
\begin{equation}
 \mathrm{ECE}=\sum_{m=1}^{M}\frac{|B_m|}{n}\left|\mathrm{acc}(B_m)-\mathrm{conf}(B_m)\right|.
\end{equation}

\item Coeficiente de correlación de Matthews (MCC): medida multiclase calculada con $s$ observaciones, $c$ predicciones correctas, $p_k$ predicciones de la clase $k$ y $t_k$ casos reales de esa clase.
\begin{equation}
 \mathrm{MCC}=\frac{c\cdot s-\sum_{k=1}^{K}p_k t_k}{\sqrt{\left(s^2-\sum_{k=1}^{K}p_k^2\right)\left(s^2-\sum_{k=1}^{K}t_k^2\right)}}.
\end{equation}
\end{itemize}

\subsubsection{Protocolo aplicado antes de los experimentos}\label{sec:protocolo-estadistico}

Antes de ejecutar cualquier comparación confirmatoria se congelan los manifiestos (identificadores, hashes y exclusiones de cada partición), se fija la separación por \texttt{lesion\_id} y se registran de antemano el conjunto de datos, la arquitectura, el preprocesamiento, las particiones, las semillas y la regla de decisión. La Figura~\ref{fig:protocolo-evaluacion} resume la secuencia de adaptación, ajuste supervisado e inferencia, incluida la separación entre la selección de variantes en el conjunto interno y la evaluación externa reservada en PAD-UFES-20 \texttt{target\_test}.

\begin{figure}[!htbp]
\caption{Protocolo de clasificación, adaptación de dominio y evaluación externa.}
\label{fig:protocolo-evaluacion}
\centering
\includesvg[width=\textwidth]{images/organizadas/21_protocolo_vertical}

\par\smallskip
\begingroup\footnotesize\singlespacing\raggedright
\par\smallskip{\fontsize{11}{13}\selectfont\singlespacing\raggedright
\noindent\textit{Nota.} Elaboración propia. El esquema se centra en el protocolo de clasificación y adaptación: D2 continúa SimSiam con el pool fuente; D3 añade \texttt{target\_adapt} a ese pool. HAM10000 aporta el ajuste supervisado y la validación interna. Cada condición ajusta su propio escalador y regresión logística; la referencia utiliza el encoder público sin adaptación. \texttt{target\_test} solo interviene en la inferencia y evaluación, con separación por paciente y sin reajuste. El esquema se apoya en DINOv2 y SimSiam \parencite{oquab2024dinov2,chen2021simsiam}. Los procedimientos de segmentación y tono se presentan en sus secciones específicas.\par}
\par\endgroup
\end{figure}

 Las comparaciones pareadas usan remuestreo bootstrap por lesión ($B=1\,000$ remuestras) con intervalo de confianza al 95\,\% por el método percentil, y corrección de Holm para contrastes múltiples: ordenados los valores $p_{(1)}\leq\cdots\leq p_{(m)}$, se rechaza la hipótesis asociada a $p_{(i)}$ si $p_{(i)}\leq \alpha/(m-i+1)$, deteniendo el procedimiento en el primer índice donde la condición falla.

Este protocolo responde directamente a cuatro riesgos de integridad de la evidencia identificados para esta línea:
\begin{itemize}
 \item \textbf{Solapamiento de imágenes de una lesión entre particiones}, que produciría fuga de información y métricas optimistas: se controla separando por \texttt{lesion\_id} y congelando listas y hashes antes de entrenar.
 \item \textbf{Bootstrap por imagen en HAM10000}, que produciría intervalos demasiado estrechos por correlación intralesión: se reporta esta limitación y se complementa con un análisis de sensibilidad agrupado por lesión antes de afirmaciones confirmatorias.
 \item \textbf{Conjuntos de datos y modalidades diferentes} (HAM10000 dermatoscópico frente a PAD-UFES-20 clínico), que podrían simular una transferencia aparente o un deterioro fuera de dominio: se describe la modalidad, la fuente y la población de cada conjunto, y PAD-UFES-20 se reporta por separado.
 \item \textbf{Línea base supervisada con partición distinta}, que haría la comparación no pareada y no atribuible solo al régimen de aprendizaje: se iguala el conjunto de datos, la arquitectura, la partición, el presupuesto de cómputo y las métricas antes de concluir.
\end{itemize}

